\chapter{3/2波动率模型设定}
\label{chap:model}

\section{3/2模型的演进与经济学直觉}
\label{sec:model_setup}

3/2模型作为一种高级的随机波动率框架，是平方根过程 (Square-Root Process，由Cox-Ingersoll-Ross普及，后被Steven Heston进一步发展) 的复杂演进形式 。

在当前的金融业界，Heston模型（即1/2模型）长期以来被视为行业标准 。然而，在面对极端的市场压力时，Heston模型往往无法准确捕捉市场中观察到的“波动率的波动率” (volatility of volatility)，并且难以完美拟合短期波动率偏斜 (short-term volatility skew) 的陡峭程度 。

为了克服这一局限性，研究人员在20世纪90年代末和21世纪初提出了一种全新的思路：将扩散项的指数从1/2修改为3/2 。这种设定的改变允许模型产生更具爆发力的波动率动态以及非线性的均值回归速度。因此，3/2模型允许波动率出现尖峰，并在市场压力下呈现出更陡峭的偏斜，这与真实的实证数据高度吻合 。然而，需要指出的是，在某些正常的市场条件下，3/2模型可能会缺乏Heston模型所具备的稳定性 。

\section{随机微分方程 (SDE) 设定}

在风险中性测度下，3/2模型的资产价格过程和瞬时方差/波动率过程可以由以下随机微分方程组来刻画 ：

标的资产价格过程 $S_t$ 的动态变化满足：

$$dS_t = r S_t dt + S_t \sqrt{V_t} (\rho dW_t^1 + \sqrt{1-\rho^2} dW_t^2)$$

其中，$r$ 代表无风险利率，$\rho$ 是两个布朗运动之间的相关系数 。同时，方差过程 $V_t$ 的动态演进满足：

$$dV_t = \kappa V_t(\theta - V_t)dt + \epsilon V_t^{3/2}dW_t^1$$

在上述方程中 ：

\begin{itemize}
    \item $\kappa$ 控制着均值回归的速度。
    \item $\theta$ 代表长期的平均方差水平。
    \item $\epsilon$ 为波动率的波动率参数，控制着方差过程的波动程度。
    \item $W_t^1$ 和 $W_t^2$ 是两个独立的标准布朗运动，分别驱动着方差过程和价格过程中的随机性。
\end{itemize}

值得注意的是，该方差过程也可以等价地表示为一种关于 $dV_t/V_t$ 的形式，从而更直观地体现其相对变化的漂移与扩散特征 。

\section{积分方差与条件分布}

在3/2模型的期权定价框架中，时间区间 $[0, T]$ 内的积分方差 (Integrated variance) $V_T$ 扮演着至关重要的角色 。其定义如下：

$$V_T = \int_0^T \sigma_t^2 dt = \int_0^T V_t dt$$

由于资产价格的随机性部分由波动率驱动，当我们以积分方差 $V_T$ 为条件时，标的资产的终端价格 $S_T$ 服从对数正态分布 (lognormal distribution) 。这一重要性质意味着，在给定 $V_T$ 的条件下，我们可以借用 Black-Scholes (BS) 模型的定价框架 。

根据伊藤等距同构 (Itô’s isometry) 原理 ：

$$\int_0^T \rho^* \sigma_t dX_t \sim \rho^* \sqrt{V_T} X_1 \sim N(0, (\rho^*)^2 V_T)$$

据此，在时间 $0$ 到 $T$ 之间的有效波动率 (effective volatility) $\sigma$ 可以表示为 ：

$$\sigma = \rho^* \sqrt{V_T / T}$$

经过进一步的随机微积分推导，终端价格与初始价格的对数收益率公式可展开为 ：

$$\log\left(\frac{S_T}{S_0}\right) = \frac{\rho}{\epsilon}\log\left(\frac{V_T}{V_0}\right) + \left(\kappa + \frac{\epsilon^2}{2}\right)V_T - \kappa\theta T - \frac{1}{2}V_T + \rho^* \sqrt{V_T} X_1$$

该公式是后续进行精确模拟和傅里叶变换定价的核心基础。